% Phiếu nhận xét của GVHD / GVPB, theo mẫu. \phieunhanxet{HƯỚNG DẪN}{hướng dẫn}{tên GV, để trống thì in dấu chấm}
\newcommand{\phieunhanxet}[3]{%
  \clearpage\thispagestyle{empty}
  {\setstretch{1.2}\fontsize{12pt}{14pt}\selectfont\noindent
  \begin{tabular}{@{}p{0.4\linewidth}@{}p{0.6\linewidth}@{}}
    \centering\bfseries \truongngan\\ \centering\bfseries KHOA CNTT\\ \centering *******
    & \centering\bfseries CỘNG HOÀ XÃ HỘI CHỦ NGHĨA VIỆT NAM\\ \centering\bfseries Độc lập -- Tự do -- Hạnh phúc\\ \centering *******
  \end{tabular}\par}
  \vspace{1cm}
  {\centering\fontsize{14pt}{18pt}\selectfont\bfseries PHIẾU NHẬN XÉT CỦA GIẢNG VIÊN #1\par}
  \vspace{0.6cm}
  \noindent Họ và tên Sinh viên 1: \sinhvienmot \hfill MSSV 1: \mssvmot\par
  \noindent Họ và tên Sinh viên 2: \sinhvienhai \hfill MSSV 2: \mssvhai\par
  \noindent Ngành: \nganh\par
  \noindent Tên đề tài: \tendetai\par
  \noindent Họ và tên Giảng viên #2: \ifstrempty{#3}{\dotfill}{#3}\par
  \vspace{0.3cm}
  {\centering\bfseries NHẬN XÉT\par}
  \begin{enumerate}[label=\arabic*.,leftmargin=*]
    \item Về nội dung đề tài \& khối lượng thực hiện:\\ \dotfill\\ \dotfill\\ \dotfill
    \item Ưu điểm:\\ \dotfill\\ \dotfill
    \item Khuyết điểm:\\ \dotfill\\ \dotfill
    \item Đề nghị cho bảo vệ hay không?\\ \dotfill
    \item Đánh giá loại: \dotfill
    \item Điểm: \dotfill
  \end{enumerate}
  \vspace{0.5cm}
  \hfill\begin{minipage}{0.6\linewidth}\centering\setstretch{1.2}
    \emph{Tp. Hồ Chí Minh, ngày \quad tháng \quad năm \namthuchien}\\
    \textbf{Giảng viên #2}\\ (Ký \& ghi rõ họ tên)
  \end{minipage}
  \clearpage}
